\chapter{社会、伦理与专业素养}

\section{领域导言与演进背景}

\begin{preamblebox}[知识领域概述与时代演进]
虽然技术议题始终是计算机课程体系的核心，但在更广阔的社会与人类文明语境中，仅有纯技术教学并不能构成一个完整的计算机高等教育体系。学生必须深入理解计算技术所植根的宏观社会背景，从而深刻把握相关的社会影响、伦理规范、法律法规及职业责任。将这些“非技术”议题系统性纳入计算机教学大纲的迫切需求，早在 1991 年 ACM/IEEE-CS 联合发布的《CS1991》课程指南中就已获得正式确认：
\begin{quote}
\small\songti
“本科生必须理解计算机学科所固有的基本文化、社会、法律与伦理问题。他们应当理解该学科的过去、现状与未来走向；理解个人在这一历史进程中的角色；并领会那些在学科发展中起到关键作用的哲学思辨、技术挑战与审美价值。

学生还需要培养针对计算技术的社会影响提出严肃质询、并批判性评估可能解答的能力。未来的从业者必须能够前瞻性地预见将某项产品引入特定环境所带来的冲击：该产品是会提升还是降低人类的生活质量？它对个人、群体乃至社会制度将产生何种深远影响？

最后，学生需要了解软硬件供应商和用户的基本法定权利，并理解构成这些权利基石的伦理价值观。未来的从业者必须清醒认识到自己肩负的社会责任，以及技术失败可能招致的严重后果。他们必须认识到自身的局限性以及所用工具的边界。所有从业者都必须做出长期承诺，在所选专业方向和整个计算机学科中保持终身学习与前沿跟进。”\footnote{引自 ACM/IEEE-CS Joint Curriculum Task Force, \textit{Computing Curricula 1991} (CS1991), ACM Press \& IEEE Computer Society Press.}
\end{quote}

随着现代信息技术的跨越式突破深刻重塑人类的生活、生产与社会治理模式，社会与伦理议题及专业职业实践的重要性和深远后果日益凸显；新型计算平台、大模型及自主决策系统每年都在提出更具挑战性的严峻难题。正是当下的学生，未来将带着审慎与担当步入工业界与学术界，承担起识别与化解这些复杂难题的崇高使命。

\textbf{课程实施路径与教学法建议：}\\
计算机教育界在讲授本领域内容时通常采取多元化模式：独立开课、在传统技术与理论课程中深度渗透、设立高阶专属研讨课、二者结合的混合模式，或作为毕业设计、工程实训和职业素养实践课程的专项模块。

本指导委员会强烈推荐\textbf{“一门必修独立核心课程 + 全程贯穿技术课程渗透”}的复合开课方案：
\begin{itemize}
    \item 一方面，本领域被列为“计算机核心”的诸多基础主题（如伦理分析哲学框架、知识产权法理、反垄断与科技经济学等）通常很难直接塞入传统算法、操作系统或编译原理等硬核技术课中。缺乏一门独立的专业核心课程，极难保证学生获得系统、完整和严谨的思辨训练。
    \item 另一方面，如果社会、伦理与职业准则仅局限于一门独立课程讲授，而在后续技术课程中绝口不提，将强化一种极度危险且错误的认知假象——即“纯粹的技术实现与伦理考量毫无瓜葛”。
\end{itemize}
因此，必须在软件工程、数据库、计算机图形学、计算机网络、系统安全、人工智能及计算机导论等多门传统技术骨干课中，以真实工程案例分析的形式，自然融入与该课程紧密相关的伦理、法律与社会思辨。例如：
\begin{itemize}
    \item 在数据库课程中探讨海量数据聚合、用户画像与数据挖掘带来的隐私侵蚀与知情同意边界；
    \item 在软件工程课程中剖析对商业客户的交付履约承诺与对终端用户、公共利益福祉之间可能产生的尖锐利益冲突；
    \item 在编程大作业中引入诸如“由非专业医师操作的高精度激光眼科手术机器人运动控制算法”等高可靠性案例，直观展示代码失误可能造成的生命伦理代价。
\end{itemize}
若专业教师对应用伦理学的教学内容与方法感到生疏，强烈建议充分利用 ACM、IEEE-CS、ACM 计算机与社会专业委员会（SIGCAS）等权威机构提供的丰富公开教学资源与案例库。

\textbf{伦理分析是贯穿所有分支的灵魂基准：}\\
伦理分析的实践构成了“社会、伦理与专业素养 (SEP)”知识领域每一个分支的底层骨架。《ACM 伦理与职业行为准则》（ACM Code of Ethics and Professional Conduct）\footnote{\url{https://www.acm.org/about/code-of-ethics}}、《IEEE 伦理准则》（IEEE Code of Ethics）\footnote{\url{https://www.ieee.org/about/corporate/governance/p7-8.html}}与《AAAI 伦理与职业行为准则》（AAAI Code of Ethics and Professional Conduct）\footnote{\url{https://aaai.org/Conferences/code-of-ethics-and-conduct.php}}共同提供了指导职业工程实践的行动规范与典型判例。其中的“一般道德律令”明确了计算机专业人员对个人责任、职业行为操守以及引领技术向善的责任基准。计算机专业教师应注重启发学生在计算机具体技术情境中的专业思辨，同时积极借助人文社科领域的优质通识课程形成学科交叉合力。

\textbf{自 CS2013 以来的核心重大变革：}
\begin{enumerate}
    \item \textbf{重磅增设“多元、公平与包容”独立核心知识单元}：自 2013 年以来，国际计算学界对正义、公平、多元和包容的认知与行动发生了深刻变革。计算机专业不仅必须为所有学生打造公平、多元与包容的受教育环境，更要求未来的计算机专业人才在投身工程实践时，确保其研发的软件、算法与系统以公平、公正与包容的方式造福全人类。为此，CS2023 正式将其升格为全领域的独立必修核心知识单元。
    \item \textbf{全面直面生成式 AI 与算法决策的社会伦理挑战}：系统涵盖大语言模型、生物特征识别、深度伪造、自动化行政裁量系统、预测性警务等新型技术在公民权利、偏见放大、劳动力替代和认知操纵方面的社会伦理冲击。
    \item \textbf{强调计算可持续性与绿色计算的紧迫性}：将硬件制造资源消耗、关键矿物提取、数据中心巨大能耗、区块链共识开销及全球电子废弃物危机纳入工程决策核心考量。
\end{enumerate}
\end{preamblebox}

\section{核心学时分配建议}

本知识领域建议的核心学时分配如下表所示。根据“向日葵模型”，全领域共涵盖 \textbf{17 个计算机核心学时}与 \textbf{14 个领域核心学时}，累计基准学时为 31 学时：

\begin{table}[H]
\centering
\setlength{\tabcolsep}{3.5pt}\fontsize{8pt}{9.6pt}\selectfont\begin{tabularx}{\columnwidth}{Y c c}
\toprule
\textbf{知识单元} & \textbf{计算机核心} & \textbf{领域核心} \\
\midrule
社会语境 & 3 & 2 \\
伦理分析方法 & 2 & 1 \\
职业伦理与专业准则 & 2 & 2 \\
知识产权与开源生态 & 1 & 1 \\
隐私与公民权利 & 2 & 1 \\
专业沟通与团队协作 & 2 & 1 \\
可持续性与绿色计算 & 1 & 1 \\
计算机发展简史 & 0 & 1 \\
计算经济学与产业生态 & 0 & 1 \\
安全策略、法律与网络犯罪 & 2 & 1 \\
多元、公平与包容 & 2 & 2 \\
\midrule
\textbf{合计} & \textbf{17 学时} & \textbf{14 学时} \\
\bottomrule
\end{tabularx}
\end{table}

\section{知识单元与详细大纲}

\subsection{社会语境}
\textbf{【学时建议】}：3 计算机核心学时，2 领域核心学时。

计算机与互联网技术在过去数十年间以超越以往任何工业发明的速度深刻重塑了整个人类社会：人类劳动生产率实现跃升，新闻资讯、大众娱乐与社交沟通渠道全面爆发，几乎在所有自然科学与工程技术领域均催生出基础性突破。社会语境单元构成了理解 SEP 领域其他所有知识单元（尤其是职业伦理）的认知根基。

\subsubsection{计算机核心主题}
\begin{itemize}
    \item \textbf{超网络化与人工智能高速演进世界的社会影响}：在人机高度互联、生成式 AI 与自主智能体加速演进的社会形态下，计算技术对劳动生产、组织协作与人类生存方式的根本性重塑。
    \item \textbf{社交媒体与人工智能对个体与意识形态的冲击}：算法推荐机制与生成内容对个体心理福祉、政治意识形态极化、信息茧房及文化认同的复杂塑造。
    \item \textbf{前沿计算技术融入公民生活与公共治理的影响}：人工智能、生物特征识别技术与算法决策系统在公民生活与行政监管中的应用剖析（如人脸识别系统、生物特征追踪标签、社会公共资源配给算法、预测性警务软件等）。
\end{itemize}

\subsubsection{领域核心主题}
\begin{itemize}
    \item \textbf{互联网、算力基础设施与人工智能的发展与控制权}：大型科技垄断巨头、跨国数据基础设施平台对公共话语场及技术治理规则的控制。
    \item \textbf{数字鸿沟及其社会外溢效应}：数字基础设施与计算资源获取能力在性别、阶层、种族、地理分布及发展中国家之间的显著失衡及其深远社会后果。
    \item \textbf{信息无障碍与暗黑设计模式}：无障碍技术的法律强制规范；以及利用心理认知陷阱诱导消费、窃取数据或阻碍注销的暗黑模式。
    \item \textbf{情境感知计算}：普适计算中设备基于环境与行为自动推断上下文所引发的社会介入与道德界限。
\end{itemize}

\subsubsection{示范性学习成果}
\begin{enumerate}
    \item \textbf{[解释]} 阐明计算机技术（计算机网络、移动计算、云计算平台）在个人交往与社会群体层面上调控和媒介人类互动的多重机制。
    \item \textbf{[解释]} 识别硬件与软件系统设计中潜藏的开发者价值预设与先验偏见，重点剖析其对边缘化群体和残障人士可用性造成的系统性障碍。
    \item \textbf{[评估]} 深度解读并阐释具体技术方案与其软硬件实现的社会语境关联。
    \item \textbf{[评估]} 运用实证经验数据，定量与定性评估特定技术设计与系统实现的社会有效性及副作用。
    \item \textbf{[解释]} 深刻阐明社交媒体算法对不同身份族群、文化群体与社会社群造成的结构性影响。
    \item \textbf{[解释]} 深入研讨互联网在促进公民、政府机构及社会各主体间沟通与互动中的赋能机制与异化风险。
    \item \textbf{[评估]} 深入剖析民主制度运行（如公共社会福利发放、电子投票系统）对计算机基础设施的依赖所带来的机遇与风险。
    \item \textbf{[解释]} 阐述历史上弱势少数族裔与女性在计算机专业领域代表性严重不足对行业文化和产品多元性造成的负面影响。
    \item \textbf{[解释]} 解释泛在普适计算系统中情境感知特性的社会伦理意涵。
    \item \textbf{[评估]} 评估不同社会形态下获取互联网和计算技术的权利鸿沟；论证为何以及如何在国际法理上将“互联网访问权”视作一项基本人权。
\end{enumerate}

\subsection{伦理分析方法}
\textbf{【学时建议】}：2 计算机核心学时，1 领域核心学时。

伦理学理论与道德哲学原则是进行理性伦理分析的基石，为复杂决策困境中的技术人员提供指引路径与思辨视角。不同的伦理学流派基于不同的哲学假设和论证方法来判定特定行为的合道德性。学生必须清醒认识到：在多元化的现实语境中，解决不同冲突可能需要借助不同的伦理学框架方能获得最佳方案；而何谓“最佳”，往往由特定语境与文化视角共同决定。应用伦理分析方法既要求深刻理解各分析工具背后的哲学原理与预设，也要求敏锐洞察决策所处的具体社会语境。源自西方哲学的传统伦理框架固然实用，但绝非放之四海而皆准；教育中必须积极引入去殖民化、本土原住民与历史上被边缘化群体的伦理视角。

\subsubsection{计算机核心主题}
\begin{itemize}
    \item \textbf{论证逻辑与逻辑谬误规避}：识别并规避技术与道德辩论中常见的非形式逻辑谬误（如偷换概念、诉诸人身、虚假对立、稻草人谬误）与误导性陈述。
    \item \textbf{伦理学理论与技术决策框架}：功利主义（后果论/效益主义）、义务论（康德道义论）、美德伦理学、关怀伦理学及社会契约论在计算技术决策中的推演与应用。
    \item \textbf{文化在技术认知中的决定性作用}：深刻理解文化背景如何塑造人们对计算技术的理解、接纳程度、产品设计与日常使用行为。
\end{itemize}

\subsubsection{领域核心主题}
\begin{itemize}
    \item \textbf{职业伦理自查清单}：工业界工程发布与算法审查中的伦理核查清单设计。
    \item \textbf{多维度评估量规}：技术方案社会效益与道德风险的量化与定性综合评估量规。
    \item \textbf{利益相关者分析}：全链路识别直接、间接利益相关者，剖析权责不对等与潜在受害者。
    \item \textbf{立场理论}：从被边缘化或受冲击最严重的弱势受众立场审视技术系统的内在缺陷。
\end{itemize}

\subsubsection{示范性学习成果}
\begin{enumerate}
    \item \textbf{[解释]} 识别并系统阐述特定文化语境与价值体系对技术工程决策产生的深刻制约与塑造。
    \item \textbf{[应用]} 熟练运用典型工程判例与哲学类比，开展严谨的技术伦理推导与论辩。
    \item \textbf{[应用]} 批判性剖析技术报告与公众论辩中的论述，准确识别并剔除常见的基础逻辑谬误。
    \item \textbf{[应用]} 严密解析复杂技术伦理争议，精准提取论辩的核心前提假设与推论逻辑链条。
    \item \textbf{[评估]} 全面评估在特定情境中所有利益相关者立场与其所处文化背景的复杂交互关系。
    \item \textbf{[开发]} 评估并量化技术决策中引入或积累“伦理债务”（Ethical Debt，即为了追求短期快速交付而推迟对伦理影响进行严谨评估所积攒的系统性潜在风险）的危害与应对策略。
\end{enumerate}

\subsection{职业伦理与专业准则}
\textbf{【学时建议】}：2 计算机核心学时，2 领域核心学时。

计算机职业伦理是应用哲学的一个重要分支，专门研究计算机从业者在专业技能与社会行为中应当遵循的决策规范。个人的职业决策主要受三重力量影响：1) 从业者个人的道德自律准则；2) 工作单位、行业协会或职业认证制度所要求的显性/隐性行为规范；3) 国际计算专业学会制定的正式职业道德与行为守则。

\subsubsection{计算机核心主题}
\begin{itemize}
    \item \textbf{社群公共价值与法律规制体系}：社会主流伦理准则、公共利益与现行法律法规的互动与统一。
    \item \textbf{专业职业素养的内涵实质}：专业能力、审慎尽责、受托人信托责任与学术/技术传承指导。
    \item \textbf{终身专业进阶与技能更迭}：保持对新兴开发工具、技术框架、法律规制与行业标准的动态跟进，具备批判性自我评估与持续职业学习能力。
    \item \textbf{专业协会职业道德准则与行为规范}：深入研习 ACM/IEEE-CS《软件工程职业道德准则和实践规范》（SE Code）、AITP、IFIP 及各国计算机学会的准则条款。
    \item \textbf{问责机制、责任归属与法律赔偿责任}：软件正确性、高可靠性与系统安全性的法律责任认定；网络安全工程师对数据隐私和防御失职的伦理保密边界。
    \item \textbf{技术哲学与社会学理论导论}：理解人类创造与使用技术的经典理论视角，包括工具主义、技术系统社会学、无障碍平权与障碍正义、技术中立论假说、实用主义、功利主义及去殖民化批判理论。
    \item \textbf{违规与道德风险的识别与举报机制}：制定策略与规范流程，敏锐识别并依法合规向上级或公众报告违反法律或职业道德准则的系统设计、软件漏洞或职业渎职行为。
\end{itemize}

\subsubsection{领域核心主题}
\begin{itemize}
    \item \textbf{计算从业人员在公共政策制定中的角色}：参与国家技术立法、标准制订、公众科普与反垄断咨询。
    \item \textbf{技术不良后果的动态前瞻与预警}：建立针对算法和系统次生灾害的持续监控与警觉机制。
    \item \textbf{伦理异议与“吹哨人”保护}：技术异议的合法提交通道、道德良知冲突下的合法检举制度及从业人员法律保护。
    \item \textbf{区域文化与伦理困境的冲突与调适}：跨国跨文化软件开发中遭遇的道德标准冲突（如言论边界、隐私感知差异）。
    \item \textbf{职场骚扰与歧视应对机制}：应对工作场所性别歧视、种族微歧视与职场霸凌的合规干预与维权途径。
    \item \textbf{职业认证与执业资质制度}：注册工程师与各专业认证的作用与局限。
    \item \textbf{工作场所计算资源的可接受使用策略 (AUP)}：雇主监控、个人隐私与企业资产安全防范边界。
    \item \textbf{人机工效学与健康计算环境}：职业性肌骨劳损、视力损害、心理重压与工效学健康标准。
    \item \textbf{上市时间、成本控制与工程质量标准的冲突}：面对商业“极速上线交付”与严格软件工程质量及安全标准冲突时的职业坚守。
\end{itemize}

\subsubsection{示范性学习成果}
\begin{enumerate}
    \item \textbf{[解释]} 敏锐识别软件需求分析、架构设计、编码实施与生产部署全流程中潜在的伦理陷阱。
    \item \textbf{[评估]} 综合权衡多方制约因素，制定妥善化解工程实践中复杂伦理困境的切实方案。
    \item \textbf{[解释]} 阐释确保软件正确性、高可靠性与人身系统安全背后的伦理与法律责任来源（如 ACM/IEEE 准则、行业强制标准、组织安全策略等）。
    \item \textbf{[解释]} 梳理专业从业者在职业生涯中持续跟踪道德准则、法律动态与最佳实践的机制与途径。
    \item \textbf{[评估]} 客观剖析相关职业道德准则的现实效力与内在局限性，评价其作为职业决策行动指南的优势与不足。
    \item \textbf{[评估]} 深入分析一个全球性计算伦理事件（如关键基础设施崩溃、重大数据泄露或波音 737 MAX 软件故障），系统剖析工程技术人员与监管官员在处置应对中的权责失衡。
    \item \textbf{[解释]} 阐明软件系统创造背后的哲学基础，及其如何决定我们选择以何种方式应用这些软件。
    \item \textbf{[解释]} 归纳计算专业从业者向政府公共政策建言献策的途径与有效形式。
    \item \textbf{[解释]} 阐明违反职业操守（如数据造假、蓄意留后门、知识产权侵权）对个人职业生涯、企业公信力及公共利益造成的灾难性后果。
    \item \textbf{[解释]} 熟知“吹哨人”道德检举机制，掌握在遭遇重大公共危害隐患时进行合规内部申诉与外部合法举报的行动指南。
    \item \textbf{[解释]} 列举并剖析区域文化差异与跨国软件工程伦理困境交互影响的真实工业案例。
    \item \textbf{[解释]} 探讨职场骚扰与歧视的各种隐蔽表现形式，并列举受害者获得组织援助与法律救济的合规通道。
    \item \textbf{[评估]} 考察国内外各类专业工程师执业资格认证制度在保障工程公众安全方面的实际效能。
    \item \textbf{[解释]} 阐释工作场所人机工效学设计与从业人员生理、心理健康之间的严密内在联系。
    \item \textbf{[解释]} 剖析工业界迫于“产品极速抢跑上市”商业压力，而妥协退让甚至践踏专业工程质量标准的严重弊端。
\end{enumerate}

\subsection{知识产权与开源生态}
\textbf{【学时建议】}：1 计算机核心学时，1 领域核心学时。

知识产权 (IP)是指人类在智力劳动成果（如计算机软件程序、算法创新、电路芯片设计等）上依法享有的无形专有权利体系。各国法律依据知识产权的具体客体属性提供不同模式的法定排他保护。在现代软件工程中，主要涉及四种截然不同的知识产权类型：专利权、版权/著作权、商业秘密以及商标权。每种权利所保障的法益对象、保护期限与排他效力各不相同。

\subsubsection{计算机核心主题}
\begin{itemize}
    \item \textbf{知识产权的哲学与法理基石}：洛克劳动财产权学说、人格权理论与功利主义激励理论对智力垄断与社会共享平衡的论证。
    \item \textbf{知识产权排他权利体系}：专有使用权、许可转让权、署名权、修改权及作品完整权保护。
    \item \textbf{无形数字知识产权 (IDIP)}：纯软件代码、数字化媒体资产、专有算法逻辑与虚拟资产的无形资产确权特征。
    \item \textbf{知识产权保护的法定框架与国际条约}：伯尔尼公约、TRIPS 协定、DMCA 法案以及各国软件著作权登记与专利审查标准。
\end{itemize}

\subsubsection{领域核心主题}
\begin{itemize}
    \item \textbf{数字版权管理 (DRM)}：技术保护措施（TPM）、加密水印与反规避条款的法理与技术博弈。
    \item \textbf{四项核心权利的边界与对比}：软件专利、软件著作权、商业秘密（如专有推荐算法源码）与商标的区别与联动防御策略。
    \item \textbf{学术与技术剽窃}：代码抄袭判定机制、开源合规审查与学术不端防范。
    \item \textbf{开源运动与开源许可协议族}：自由软件基金会（FSF）与开源倡议组织（OSI）理念；宽松型许可证（MIT、Apache 2.0、BSD）与传染型许可证（GPLv2/v3、AGPL）的法律约束与商业合规。
    \item \textbf{软件盗版与网络侵权}：全球软件盗版的经济成因、暗网盗版分发链条与反侵权救济。
\end{itemize}

\subsubsection{示范性学习成果}
\begin{enumerate}
    \item \textbf{[解释]} 结合具体国家或区域法域，深入辨析知识产权保护的哲学伦理基础（激励创新 vs 知识垄断弊端）。
    \item \textbf{[解释]} 系统阐明旨在遏制数字版权侵权行为的现行国际与国内法律法规。
    \item \textbf{[评估]} 批判性审视和评价数字版权相关法案在防范侵权与保障公众合理使用之间的失衡争议。
    \item \textbf{[应用]} 准确识别当今数字化前沿场景中涌现的各种无形数字知识产权形态（如 AI 生成物、数据集所有权、NFT）。
    \item \textbf{[评估]} 在具体工程研发情境下，合法合理论证对受版权保护素材进行“合理使用”的法律抗辩依据。
    \item \textbf{[评估]} 评估学术界与工业界各类自动化代码抄袭检测机制所引发的技术误判与潜在隐私伦理争议。
    \item \textbf{[应用]} 精准解读主流开源许可协议（如 GPL、Apache、MIT）的核心条款约束，评估其在商用产品集成中的合规风险。
    \item \textbf{[评估]} 全面权衡为软件算法申请专利保护与作为商业秘密保护之间的利弊得失与法律壁垒。
    \item \textbf{[解释]} 严谨区分并对比著作权、专利权、商业秘密和商标在保护计算机软件产品时的法律维度差异。
    \item \textbf{[解释]} 阐明开源运动的发展历史、核心价值观与推动现代软件基础设施繁荣的战略意义。
    \item \textbf{[解释]} 剖析全球范围内软件盗版与逆向工程的跨国治理困境与产业应对措施。
\end{enumerate}

\subsection{隐私与公民权利}
\textbf{【学时建议】}：2 计算机核心学时，1 领域核心学时。

电子信息的低成本存储、秒级跨网传播与海量聚合，使得在“保护个人隐私权益”与“促进信息自由流通和公共治理利用”之间寻求精巧平衡成为极其紧迫的时代课题。数字化数据采集的极度便捷，使个人隐私权与公民自由面临前所未有的制度性侵蚀，而不同文化传统与主权国家对隐私边界的定义与敏感度存在巨大差异。

\subsubsection{计算机核心主题}
\begin{itemize}
    \item \textbf{泛在数据采集的隐私侵蚀}：日常交易数据库、超大规模数据仓库、公共监控摄录系统、车联网传感器与云计算对个人全生命周期数字足迹的无死角记录。
    \item \textbf{差分隐私（Differential Privacy）数学原理与应用外溢}：隐私预算参数 $\varepsilon$ 的权衡、拉普拉斯噪声注入及在人口普查与用户画像聚合中的实践效应。
    \item \textbf{隐私保护技术增强方案}：同态加密、零知识证明、安全多方计算（SMPC）、端到端加密通信与联邦学习。
    \item \textbf{公民权利、言论空间与跨文化法理差异}：不同国家与文化背景下对信息知情权、被遗忘权、匿名权与公民言论自由的法理界定与实践冲突。
\end{itemize}

\subsubsection{领域核心主题}
\begin{itemize}
    \item \textbf{个人隐私权的哲学与宪法法理渊源}：隐私作为人格尊严核心要素的哲学论证，隐私权作为消极自由与自主控制权的法律演进。
    \item \textbf{重点法域的隐私保护立法框架}：欧盟通用数据保护条例（GDPR）、加州消费者隐私法案以及中国《个人信息保护法》（PIPL）的核心合规准则。
    \item \textbf{垂直专业领域的合规要求}：医疗健康数据（HIPAA）、金融交易信用数据、未成年人网络信息保护法等专项合规。
    \item \textbf{表达自由与有害言论规制的张力}：言论自由保护与遏制虚假信息、仇恨言论、网络诽谤之间的法理平衡界限。
\end{itemize}

\subsubsection{示范性学习成果}
\begin{enumerate}
    \item \textbf{[评估]} 针对高并发交易数据库与海量数据仓库系统中的数据泄露威胁与越权访问漏洞，提出完备的技术防御与审计方案。
    \item \textbf{[解释]} 系统阐明现代泛在监视系统（如 RFID 追踪、智慧城市人脸识别天网、ETC 轨迹追踪、智能手机地理围栏）背后的数据采集机制与社会控制隐患。
    \item \textbf{[解释]} 严密阐述差分隐私（Differential Privacy）的核心数学原理、隐私损失边界及在保护统计隐私中的实际成效与局限。
    \item \textbf{[评估]} 深入调研并评估新兴密码学方案与隐私增强技术（如联邦学习、同态加密）在工业级数据中台中的落地代价与防御成效。
    \item \textbf{[解释]} 探讨在不同司法管辖区及跨文化语境中，确立个人隐私法律保护的哲学法理渊源。
    \item \textbf{[评估]} 批判性审视各类现行隐私立法（如 GDPR、PIPL）在立法初衷、合规成本、遏制技术垄断与阻碍初创企业创新之间的多维博弈。
    \item \textbf{[应用]} 提出在设计社交内容与公共交流平台时，既保障用户正当言论自由又有效治理恶意违规内容的系统性工程策略。
\end{enumerate}

\subsection{专业沟通与团队协作}
\textbf{【学时建议】}：2 计算机核心学时，1 领域核心学时。

专业沟通的本质是向具有不同背景、专业知识储备和诉求目标的受众精准、清晰地传递复杂的技术与管理信息。高水平的专业技术沟通能力绝非与生俱来的天赋，而必须在整个大学本科课程体系中依托真实工程情境反复操练与持续雕琢。

\subsubsection{计算机核心主题}
\begin{itemize}
    \item \textbf{硬核技术资料的深度解读、提炼与综合}：深入研读复杂工业源码、架构规范文档、RFC 协议标准并精准撰写执行摘要与技术白皮书。
    \item \textbf{高质量工程技术文档编制}：编写规范严密的技术教程、开发者参考指南及生产级 OpenAPI/SDK 接口文档。
    \item \textbf{团队多模态高效沟通准则}：在口头汇报、书面备忘录、即时通讯工具与异步协作流中践行“清晰、礼貌、精炼”沟通规范。
    \item \textbf{受众感知与针对性沟通策略}：敏锐识别受众画像（商业客户、非技术管理层、资深架构师、法务合规人员），动态调整技术表达颗粒度与侧重点。
    \item \textbf{现代化协同开发工具链的规范运用}：Git/GitHub 协作流、Issue 跟踪系统、敏捷看板、知识库 Wiki 与远程协同评审工具的深度使用。
    \item \textbf{技术争议中的修辞谬误防范}：在代码评审与技术路线争辩中彻底摈弃人身攻击与情感绑架，依靠实测数据与逻辑推导化解分歧。
    \item \textbf{面向专业受众的无障碍与包容性表达}：图表可读性设计、色盲友好配色、屏幕阅读器可识别文本及无歧视中立技术词汇。
\end{itemize}

\subsubsection{领域核心主题}
\begin{itemize}
    \item \textbf{书面与口头交流中的跨文化能力}：在跨国、多语言、多元文化分布式团队中展现文化敏锐度，规避文化禁忌与语言歧义。
    \item \textbf{多维竞争价值观的权衡与精炼表达}：向高管与客户精准汇报技术指标、研发进度、系统质量、人力成本、市场窗口期与财务预算之间的妥协与权衡。
    \item \textbf{以专业写作推动职场问题解决与决策}：撰写正式建议书（RFC）、向管理层警示技术债务与伦理风险、提交针对系统易用性与无障碍改进的技术论证报告。
\end{itemize}

\subsubsection{示范性学习成果}
\begin{enumerate}
    \item \textbf{[解释]} 深刻理解遵循严格排版格式、准确绘制架构图表与规范引用权威文献对撰写专业技术文档的至高重要性。
    \item \textbf{[评估]} 从技术严谨性、表达精炼度、术语无歧义性及受众契合度等多重维度，对他人撰写的技术文档开展专业评审与同行反馈。
    \item \textbf{[开发]} 面向多元化专业受众，设计并完成一场结构严整、图文并茂、照顾无障碍需求且论点鲜明的正式技术汇报演说。
    \item \textbf{[应用]} 规划并主持团队协作活动（线上协作、线下碰头、跨时区异步文档共创），主动营造鼓励全员包容性参与、倾听并珍视不同创见的良性协作氛围。
    \item \textbf{[解释]} 归纳并践行高效团队沟通的黄金法则（明确行动项、降低沟通熵值、文档留痕与闭环反馈）。
        \item \textbf{[应用]} 探讨在多元化与跨文化工程团队中，促进团队高效协作并提升产出成果的有效方法。
    \item \textbf{[评估]} 客观评估自身在依托多元背景协同工作的远程分布式团队中的个人沟通优势与短板，提出针对性提升计划。
\end{enumerate}

\subsection{可持续性与绿色计算}
\textbf{【学时建议】}：1 计算机核心学时，1 领域核心学时。

联合国将“可持续发展”定义为：“既满足当代人的需求，又不对后代人满足其自身需求的能力构成危害的发展模式”；亦可概括为“生态环境、社会公平与经济效益之间的动态平衡”。随着计算技术全方位渗透人类社会，据权威测算，全球电力消耗中已有约 10\% 用于支撑各类计算与网络系统，且该比例仍在加速攀升。此外，消费级电子硬件的极速淘汰引发了对稀土元素与贵金属开采的畸形需求，并催生出触目惊心的全球电子废弃物生态危机。未来的计算机从业者必须牢固树立生态意识，在工程设计与架构决策中深刻预见计算基础设施给全球环境与地方社区带来的长期资源与碳足迹负担。

\subsubsection{计算机核心主题}
\begin{itemize}
    \item \textbf{践行可持续发展理念的工程实践}：在算法与系统实现中全方位权衡环境开销、社会影响、经济可行性与计算能耗。
    \item \textbf{计算设备全生命周期的生态足迹}：计算芯片制造过程中的高能耗与水资源消耗、稀土采掘污染、设备淘汰报废及跨国电子垃圾倾销对发展中国家脆弱生态的重创。
    \item \textbf{高能耗共识算法与能耗权衡}：以区块链比特币 PoW（工作量证明）算法与以太坊 PoS（权益证明）算法为典型范例，深入量化评估其计算能耗、算力军备竞赛与碳排放代价。
\end{itemize}

\subsubsection{领域核心主题}
\begin{itemize}
    \item \textbf{绿色 IT 与低碳软件设计标准规范}：绿色软件工程协会能效标准、能效评级指标（PUE）与软件碳强度（SCI）规范。
    \item \textbf{技术媒介复杂现象的系统性生态效应}：大规模居家远程办公、跨境算力外包、电商平台物流算法对整体社会能源消费结构的系统性级联影响。
    \item \textbf{普适计算与智慧环保赋能}：利用物联网智能电网、智慧交通路线规划、环境自动监测无线传感器网络、公民科学与环境激进主义系统推动社会实现低碳转型。
    \item \textbf{计算技术在解决前沿环境问题中的应用研究}：利用高性能计算与 AI 大模型攻关新能源开发、大气污染扩散模拟、碳循环捕集、循环经济资源调度与精准农业。
    \item \textbf{软件系统可持续性与社会制度的深度共生}：系统生命力与用户认知能力、企业组织流程及政府强制碳税政策之间的互赖机制。
\end{itemize}

\subsubsection{示范性学习成果}
\begin{enumerate}
    \item \textbf{[应用]} 掌握在日常软件工程与系统架构实践中践行“绿色与可持续工程师”的有效途径与衡量准则。
    \item \textbf{[评估]} 针对任意给定的计算工程项目（如大型数据中心上云、开发训练超大参数量模型、投产边缘物联网终端），系统列举并量化其全生命周期的环境与资源开销。
    \item \textbf{[解释]} 绘制并阐明全球电子废弃物跨国转移的生态危害链条，反思科技产品“计划性报废”商业模式的道德伦理失范。
    \item \textbf{[评估]} 客观列举并定量测算远程办公、在线电商或加密货币挖矿等现代数字化活动的综合能源增减效应。
    \item \textbf{[评估]} 深入分析算法时间复杂度、操作系统任务调度、分布式缓存、数据库查询优化等计算机内核设计决策对物理硬件功耗的直接影响。
    \item \textbf{[评估]} 调研全新大型软件系统架构投产对所在城市微电网、散热水资源等周边物理环境的长期负外部性。
    \item \textbf{[开发]} 依据国际绿色计算标准，编制一套适用于本企业软件研发与机房部署的绿色低碳开发操作指南。
    \item \textbf{[评估]} 调研智慧能源调度、智能物流供应链管理及环境遥感监测等计算系统在化解现实生态环境危机中的落地价值。
\end{enumerate}

\subsection{计算机发展简史}
\textbf{【学时建议】}：1 领域核心学时。

开设计算机科学史并非为了单纯记诵历史事实，而是为了让学生深刻领悟技术的高速演进如何在长周期内深刻塑造人类文明。科学史教学通常与“系统基础 (SF)”和“软件开发基础 (SDF)”紧密结合。研习历史使我们能够洞悉当今计算机系统架构为何如此设计、早期方案是在何种具体的社会需求与硬件技术瓶颈下诞生、以及历史的技术路径依赖如何持续在当今学术与产业界激荡回响。

\subsubsection{领域核心主题}
\begin{itemize}
    \item \textbf{第一纪元：史前文明与机械计算（1946 年 ENIAC 诞生前）}：古代模拟计算工具（巨石阵天文历法、安提基特拉机械仪、中世纪机械大钟）；欧几里得几何算法、阿达·洛夫莱斯与查尔斯·巴贝奇分析机、哥德尔不完备定理、阿隆佐·邱奇Lambda 演算、艾伦·图灵图灵机；机电式与继电器计算硬件。
    \item \textbf{第二纪元：现代早期数字电子计算}：ENIAC、UNIVAC、布莱切利庄园破译密码的密码机、晶体管替代电子管革命（TRADIC）、大型主机与最早的通用操作系统（IBM 操作系统/360）。
    \item \textbf{第三纪元：现代微型计算与个人电脑革命}：微处理器诞生、微型计算机（PC）普及、摩尔定律（Moore's Law）的半个世纪奇迹、施乐帕克研究中心对图形用户界面（GUI）与以太网的奠基、第一代 Apple 计算机与个人计算普及。
    \item \textbf{第四纪元：互联网时代与分布式互联}：ARPANET、TCP/IP 协议栈统一世界、万维网（WWW）诞生与浏览器大战、互联网技术标准委员会（IETF/W3C）、互联网商业泡沫与科技巨头垄断格局形成。
    \item \textbf{第五纪元：智能手机、移动互联与云计算时代}：智能手机生态普及（iOS、Android 移动革命）、集中式多租户云计算基础设施（AWS、Azure）、软件即服务、全网移动社交网络繁荣与隐私危机的全面交织。
    \item \textbf{第六纪元：AI 辅助与前沿智能时代}：深度学习复兴、基于算法的自动化决策支持系统、大规模推荐引擎、大语言模型（LLMs）与生成式人工智能对内容生产、工程研发与人机交互形态的全面重塑。
\end{itemize}

\subsubsection{示范性学习成果}
\begin{enumerate}
    \item \textbf{[解释]} 清晰辨析计算机学科发展史上几个关键历史时期的技术代际特征与核心演进驱动力。
    \item \textbf{[解释]} 准确阐述多位计算科学先驱（如图灵、冯·诺依曼、洛夫莱斯、香农、戴克斯特拉）对现代计算理论与工程架构的不可磨灭贡献。
    \item \textbf{[评估]} 深入探讨计算机发展史上重大里程碑事件的历史语境与深远意义（例如：电子管向晶体管与集成电路的过渡、操作系统/360 的软件工程危机、施乐 PARC 实验室的人机交互革命、第一款计算机病毒发端、WWW 协议开源、千年虫危机 Y2K、移动智能机浪潮）。
    \item \textbf{[评估]} 深刻对比个人电脑普及与互联网大爆炸前后，人类社会日常生活、工作协作与全球经济版图的天翻地覆变革。
\end{enumerate}

\subsection{计算经济学与产业生态}
\textbf{【学时建议】}：1 领域核心学时。

计算经济学不仅关乎软硬件服务提供商的生存壮大，也直接牵涉全社会各行业的经济运行，同时与数以亿计的专业及非专业终端用户的切实福祉紧密相连。计算机从业人员必须理解计算产业运行的经济学规律，方能在复杂的商业浪潮中做出清醒的技术选型与战略研判。

\subsubsection{领域核心主题}
\begin{itemize}
    \item \textbf{数字服务提供商的市场生态}：强监管与弱监管市场环境、自然垄断与开放竞争市场；大型科技跨国平台构建的“围墙花园”专有生态及其对跨平台创新的抑制。
    \item \textbf{知识经济与注意力经济}：数据要素资产化、算法诱导机制以争夺人类有限注意力作为核心变现模式的经济本质。
    \item \textbf{专业高技术人才供求平衡与技术产品质量的关系}：全球算力与软件人才缺口对系统架构质量、测试充分性与安全隐患的深层连锁反应。
    \item \textbf{计算领域的商业定价与商业模式演化}：买断制、软件订阅制、计划性报废（人为制造硬件换代需求）、软件使用许可授权、自由及开源软件商业化路径。
    \item \textbf{软件工程外包与离岸研发}：全球研发分工、离岸交付中心对发达国家本土就业与发展中国家技术生态的双向冲击。
    \item \textbf{全球化分工与科技自主的博弈}：跨国技术产业链分工与地缘政治科技壁垒对开发者生态和最终用户的交织影响。
    \item \textbf{计算资源获取鸿沟造成的经济不平等}：高端算力垄断与中小企业技术研发创新的两极分化。
    \item \textbf{前沿智能化对劳动力市场的颠覆性重构}：机器人自动化与生成式 AI 对白领、程序员及基础劳动岗位的结构性冲击与就业转型阵痛。
    \item \textbf{规模经济、科技创业孵化与科技慈善}：零边际复制成本、高固定研发投入特征下的风险投资（VC）模式与开源基金会慈善运作。
    \item \textbf{计算技术颠覆现代个人与社会金融体系}：分布式区块链与加密货币、移动支付与数字钱包、发展中国家基于短信（SMS）的普惠金融基础设施变革。
\end{itemize}

\subsubsection{示范性学习成果}
\begin{enumerate}
    \item \textbf{[解释]} 归纳并剖析科技巨头垄断、平台“围墙花园”封闭生态对数字自由竞争与开放 Web 生态的扼杀隐患。
    \item \textbf{[解释]} 探讨高水平专业技术人才严重短缺对企业级核心系统稳定性、安全性保障造成的直接拖累。
    \item \textbf{[解释]} 梳理计算软件与数字服务商业模式自早期纸带/软盘买断，走向订阅制、按用量计费与免费增值的历史演进脉络。
    \item \textbf{[解释]} 深入阐释“知识经济”与“注意力经济”的深层社会代价（如短视频沉迷、心理焦虑、深度虚假信息传播）。
    \item \textbf{[评估]} 总结科技全球化分工与近期部分国家盛行的逆全球化科技壁垒对全球开源协作与技术标准统一造成的双重影响。
    \item \textbf{[评估]} 深刻分析人工智能全自动化浪潮对人类劳动力市场结构的冲击，前瞻性预判传统技术开发岗位的演化方向。
    \item \textbf{[解释]} 详细阐明计算技术基础设施的升级如何彻底重塑现代企业的组织架构与管理范式。
    \item \textbf{[评估]} 系统梳理移动支付、区块链金融与数字法币对个人理财及国家金融安全带来的积极赋能与潜在洗钱、投机金融风险。
\end{enumerate}

\subsection{安全策略、法律与网络犯罪}
\textbf{【学时建议】}：2 计算机核心学时，1 领域核心学时。

虽然安全策略、法律规范与网络犯罪惩治是现代计算体系的核心支柱，但这些议题绝非单纯的技术对抗，必须植根于“知识产权”、“隐私与公民权利”、“社会语境”及“职业伦理”等先修知识的宏观视野中审视。计算机与互联网一方面为人类文明带来空前繁荣，另一方面也孕育出前所未有的隐私威胁、新型跨国有组织计算机犯罪、反社会网络暴力、企事业单位运营瘫痪以及关键基础设施的系统性高危集中风险。

\subsubsection{计算机核心主题}
\begin{itemize}
    \item \textbf{经典计算机犯罪典型案例与法律追责机制}：重大网络勒索攻击、分布式拒绝服务（DDoS）敲诈、特大跨国数据泄露及执法机构的跨国司法管辖追诉。
    \item \textbf{社会工程学、高技术欺诈与灾后响应恢复}：网络钓鱼、深度伪造欺诈、身份冒用以及被入侵后的应急灾备恢复机制。
    \item \textbf{计算机犯罪的法律界定范畴}：非法侵入计算机信息系统罪、越权获取数据、破坏计算机信息系统功能、非法利用信息网络协助犯罪。
    \item \textbf{网络恐怖主义与黑客地下黑产攻击动机及破坏后果}：国家级高级持续性威胁（APT）、关键市政能源基础设施破坏、地下黑客暗网数据买卖黑产链。
    \item \textbf{恶意软件谱系及其系统破坏机理}：计算机病毒、自主传播蠕虫、特洛伊木马、间谍软件、勒索加密软件与 Rootkit 的行为特征。
\end{itemize}

\subsubsection{领域核心主题}
\begin{itemize}
    \item \textbf{企事业单位安全防护与犯罪预防战略}：纵深防御架构、网络攻防演练（红蓝对抗）、零信任安全接入规范与安全意识全员培训。
    \item \textbf{企业级信息安全策略制定与合规审计}：最小权限准则、强密码生命周期管理、多因素认证策略、雇员合规网络行为审计准则。
\end{itemize}

\subsubsection{示范性学习成果}
\begin{enumerate}
    \item \textbf{[解释]} 梳理引发广泛社会震荡的经典计算机犯罪事件与社会工程学攻击案，总结其背后的技术与管理脆弱点。
    \item \textbf{[解释]} 准确列举并理解适用于惩治网络黑客入侵、数据窃取与网络诈骗的相关刑事法律与网络安全法条。
    \item \textbf{[解释]} 深入剖析网络恐怖主义与恶意破解者的心理动机、行为特征以及对国家安全与社会秩序的破坏性后果。
    \item \textbf{[评估]} 深刻审查因权限滥用、内部特权账号失控所导致的严重数据泄露事件中的法律违规与职业道德背叛。
    \item \textbf{[评估]} 深入探讨计算机专业人员在构筑信息系统安全防线中的关键角色，以及在“绝对安全”、“系统可用性”与“用户隐私”之间寻求权衡的职业取舍。
    \item \textbf{[应用]} 针对企事业单位及个人，系统设计一套抵御计算机勒索犯罪、防范身份冒用及防范跨国欺诈的综合防治方案。
    \item \textbf{[开发]} 为一家虚拟中型科技企业独立编制一份完备的企业级安全策略文档，涵盖口令密码策略、特权访问分级、员工行为合规监控及数据出境保护规程。
\end{enumerate}

\subsection{多元、公平与包容}
\textbf{【学时建议】}：2 计算机核心学时，2 领域核心学时。

自计算机学科建立伊始，多元性不足就是困扰该领域的长期顽疾。尽管计算机专业充满创造力、薪资待遇丰厚且就业空间广阔，但在种族、性别、身体机能及社会背景方面的代表性严重失衡现象却普遍存在。对于极其广泛的弱势群体而言，他们的第一门计算机先修课往往不幸成为他们的最后一门课。这种失衡背后交织着系统性的排斥壁垒、刻板印象、健全主义偏见、性别偏见与阶层固化等多重历史与制度成因，绝非依靠某种单一速效方案所能根治。

CS2023 的发起主办学会——ACM、IEEE-CS 与 AAAI，均将“多元、公平与包容（DEI）”确立为协会的核心准则。\footnote{参见 ACM DEI 准则：\url{https://www.acm.org/diversity-inclusion/about}；IEEE 多元化声明：\url{https://www.ieee.org/about/diversity-index.html}；AAAI 多元化政策：\url{https://aaai.org/Organization/diversity-statement.php}}计算机高等教育课堂不仅应当向学生传授这些原则，更应成为践行与倡导包容性文化的模范阵地：
\begin{itemize}
    \item \textbf{多元}：真诚拥抱并主动吸纳人类全方位的个体差异（包括性别认同、种族族群、身体功能状况、文化信仰、国籍背景等）；
    \item \textbf{包容}：致力打破物理与数字藩篱，消除无障碍壁垒，营造让每一位成员感受到被尊重、被倾听与拥有归属感的友好学术社群；
    \item \textbf{公平}：直面历史不公，精准倾斜教育资源，审查与校正无意识偏见，确保不同背景的人才均能获得平等的成就空间。
\end{itemize}
在计算机课程体系中全面贯穿 DEI 理念不仅具有坚实的道德与正义正当性，更是决定科技产业健康发展的关键能力。权威实证数据表明，背景多元化的研发团队在创造力与抗风险能力上显著超越同质化团队；许多震惊世界的重大技术溃败事故（如人脸识别在深肤色群体上的极高误判率、早期汽车安全气囊未适配女性与儿童体征致死悲剧），原本完全可以通过构建多元视角的研发与评测团队予以彻底规避。

\subsubsection{计算机核心主题}
\begin{itemize}
    \item \textbf{个体身份与计算技术环境的双向互塑}：个人身份特质如何深刻影响在学术与产业计算环境中的职业体验，以及技术系统反向对不同群体身份认同造成的赋能或标签化禁锢。
    \item \textbf{多元化研发团队的战略收益与同质化团队的重大盲区}：深入剖析团队成员背景多元化带来的算法盲区规避、产品健壮性提升以及同质化带来的技术偏见风险。
    \item \textbf{包容性语言与敏感偏见术语的净化规范}：掌握国际计算界推动的中立包容用语规范，规避带有历史压迫色彩的争议性词汇（如 Master/Slave 架构重构为 Primary/Replica、Blacklist/Whitelist 重构为 Denylist/Allowlist）。
    \item \textbf{包容性行为范式及其职场价值}：在项目评审、日常研发与决策讨论中主动践行倾听弱势声音、提供心理安全感的包容行为。
    \item \textbf{无障碍辅助技术与通用设计}：符合国际 W3C Web 内容无障碍指南（WCAG）规范的交互系统开发、辅助屏幕阅读器适配与无障碍软硬件基准。
    \item \textbf{计算机从业人员对社会正义的正负向杠杆作用}：从业者通过其编写的代码和算法，既可能加剧社会歧视与阶层鸿沟，也拥有强大的正向杠杆来促进社会公平与正义。
\end{itemize}

\subsubsection{领域核心主题}
\begin{itemize}
    \item \textbf{树立契合现实世界的多元化榜样}：在课程与教材中充分展示具有不同性别、种族和身体状况的杰出计算先驱、行业专家与高年级优秀学长。
    \item \textbf{多元化的系统性实证收益与同质化带来的高昂代价}：深入研讨多元化提升软件工程质量与规避重大事故的经典工业界实证调研与案例剖析。
    \item \textbf{技术霸权造成的历史边缘化与全球数字基建公平挑战}：西方技术霸权语境下边缘群体的技术失语症、全球算力基建分布不均对发展中国家公平获取技术红利的阻碍。
\end{itemize}

\subsubsection{示范性学习成果}
\begin{enumerate}
    \item \textbf{[解释]} 准确界定并辨析“公平”、“平等”、“多元”与“包容”在计算机工程语境中的深刻区别。
    \item \textbf{[解释]} 深刻阐明技术权力与特权结构在计算机行业文化、软硬件产品定义及社会阶层流动中扮演的角色。
    \item \textbf{[解释]} 总结能够营造包容性职场与团队氛围的日常表达用语、工程习惯与协作行为；主动践行 ACM “Words Matter”包容性用语指南。\footnote{参见 ACM Words Matter 倡议：\url{https://www.acm.org/diversity-inclusion/words-matter}}
    \item \textbf{[评估]} 依据 W3C WAI/WCAG 国际标准，独立开展本专业实验室、教学机房及学院官方主页的信息无障碍评测，精准定位可访问性障碍并提出修复方案。
    \item \textbf{[应用]} 具备与不同身份背景、文化信仰的团队伙伴精诚合作、相互尊重的团队协同能力。专业教师绝不能仅将学生随机分组了事，而应全程指导、监控并培养学生团队的包容性协作能力。
    \item \textbf{[评估]} 采集并比对所在高校计算机及 STEM 专业的人口统计学画像与全校整体学生画像的差异；深入剖析导致特定群体在计算机专业招生、留存与毕业深造中处于劣势的制度与心理成因。
    \item \textbf{[评估]} 比对所在高校或企业的人口结构与周边社会社区的人口多样性差异，评估科技高地是否存在与本地社区脱节的“精英技术孤岛”效应。
    \item \textbf{[解释]} 敏锐识别软硬件系统设计中暗含的开发者文化预设与潜意识偏见，评估其对不同文化背景和身体障碍用户的可用性杀伤力。
    \item \textbf{[解释]} 在课程技术案例中积极引用并致敬来自不同历史时期、跨文化背景的杰出计算科学家，打破行业单一刻板印象。
    \item \textbf{[解释]} 列举生动案例，证明研发团队多元化如何帮助企业成功开拓全球新兴市场，以及因缺乏多元视角而导致重大产品翻车的教训。
    \item \textbf{[开发]} 针对熟悉的日常场景（如《程序设计基础》初级核心课），提出一套切实可行的系统性改进方案，从教学案例选择、答疑辅导机制到助教培训全面促进多样、公平与包容。
\end{enumerate}

\section{专业素养与思维倾向}

在社会、伦理与专业素养领域，单纯掌握知识点不足以确保毕业生在职业生涯中做出合乎道德的技术行动。CS2023 强调必须内化以下六大核心素养特质：
\begin{itemize}
    \item \textbf{专业职业精神}：现代计算实践极少是孤立封闭的技术操练，高度依赖团队配合、客户协同、组织领导及同行交流。无论面对何种身份特质与文化背景的协作伙伴，均能展现出高度的同理心、职业严谨性与平等待人风范。
    \item \textbf{责任担当意识}：深怀对代码质量、系统可靠性与社会后果的敬畏之心，对微小架构漏洞与潜在恶性连锁反应保持极高的责任审慎感。
    \item \textbf{批判性自我反思}：具备深度审视自身行为、潜意识偏见、所享特权及深层动机的思辨反思能力。主动学习识别隐性偏见，并在算法建模和团队决策中持续纠偏。
    \item \textbf{环境敏捷响应性}：能够迅速、审慎地适应外部技术与社会环境的突变（如面对突发公共事件时从线下集中办公平稳切换至分布式远程研发模式），并主动重构对“专业性”内涵的认知。
    \item \textbf{前瞻主动性}：在职场中不仅被动执行规则，更能主动洞察新兴前沿趋势（如国际无障碍标准升级、绿色算法优化），并前瞻性地将最佳实践推行至团队现有研发流程中。
    \item \textbf{文化能力}：高度重视跨文化协同能力，在跨国分布式开发中规范使用包容性语言，警惕文化偏见，主动促进真诚、透明与平等的跨团队思想交锋。
\end{itemize}

\section{数学基础要求}

SEP 领域虽然重点关注社会学、法学与伦理学，但其定量评估部分与现代应用数学紧密相连：
\begin{itemize}
    \item \textbf{离散数学与形式逻辑}：论证有效性验证、逻辑谬误拆解、形式化安全策略验证与知识产权权利范围的严密界定。
    \item \textbf{概率论与数理统计}：算法公平性定量测算（如人口均等 Demographic Parity、均等化赔率 Equalized Odds）、差分隐私拉普拉斯扰动注入与隐私预算 $\varepsilon$ 评估、多元化实证调研抽样推断。
\end{itemize}

\section{交叉主题与跨学科共享概念}

\subsection{与全书其他知识领域的共享概念}
\begin{itemize}
    \item \textbf{社会正义、公平、多元与包容}：贯穿全书所有 17 个 KA，尤其是人工智能（AI）、人机交互（HCI）、软件工程（SE）和计算机系统（系统基础/OS）。
    \item \textbf{超网络化与 AI 高速演进世界的社会影响 (社会伦理/社会情境与计算影响)}：与 人工智能/人工智能基础议题、人机交互/交互基础、安全隐私/Fundamental Concepts 紧密相扣。
    \item \textbf{互联网的演进与控制权治理}：与 网络通信/网络导论与体系结构 互联网体系结构与组织管理共享。
    \item \textbf{情境感知计算}：与 人机交互/Design for non-mouse interfaces 普适与情境感知计算深度交叉。
    \item \textbf{职业认证与伦理行为规范}：与 安全隐私/Fundamental Concepts 伦理与合规要求互为支撑。
    \item \textbf{知识产权与数字专有权}：与 数据管理/Information Storage and Retrieval 知识产权与数字资产保护协同。
    \item \textbf{广泛数据采集带来的隐私侵蚀}：与 数据管理/Database Systems 数据独立性、数据管理/Data Mining 数据清洗与挖掘伦理直接关联。
    \item \textbf{基于技术的隐私保护解决方案}：与 安全隐私/Threats and Attacks 针对隐私与匿名性的对抗攻击互为攻防。
    \item \textbf{隐私权的哲学渊源}：与 人工智能/人工智能基础议题 哲学伦理思辨深度交织。
    \item \textbf{高效口头、书面与团队工程沟通 (社会伦理/专业沟通与表达)}：与 人机交互/Collaboration and Communication 群体沟通、软件工程/项目管理与度量 敏捷团队协作深度重叠。
    \item \textbf{现代计算发展史}：与 体系结构/Digital logic and digital systems 计算机体系结构演进史紧密映射。
    \item \textbf{计算机犯罪、社会工程学与安全策略}：与 安全隐私/Digital Forensics 证据采信规则、人机交互/Human Factors and Security 信任与欺骗建模、安全隐私/Security Policy and Governance 企业安全治理深度协同。
\end{itemize}

\subsection{贯穿全领域的横切主题}
\begin{itemize}
    \item \textbf{正义} \quad $\bullet$ \textbf{公平} \quad $\bullet$ \textbf{多元} \quad $\bullet$ \textbf{包容}
    \item \textbf{社会议题} \quad $\bullet$ \textbf{计算伦理} \quad $\bullet$ \textbf{专业职业素养}
\end{itemize}

\section{课程体系规划与开课建议}

联合工作组建议，各高校计算机系至少应在其他技术核心课程中全面涵盖 计算机核心 的学习目标；理想状态下，领域核心 的核心内容也应全方位渗透：

\subsection{中阶独立课程方案 (面向 17 计算机核心学时)}
在本科二年级或三年级设立一门 17--20 学时的必修独立课程《计算机职业素养与社会伦理导论》，紧凑覆盖所有 计算机核心 单元：
\begin{itemize}
    \item 社会语境：3 学时
    \item 伦理分析方法：2 学时
    \item 职业伦理与专业准则：2 学时
    \item 知识产权：1 学时
    \item 隐私与公民权利：2 学时
    \item 专业沟通：2 学时
    \item 可持续性与绿色计算：1 学时
    \item 安全策略、法律与网络犯罪：2 学时
    \item 正义、公平、多元与包容 (DEI)：2 学时
\end{itemize}

\subsection{中高阶独立完整深化方案 (面向 31 全核心学时)}
在本科高年级设立一门标准的 32 学时（2 学分）完整必修课《计算伦理、法律、社会与可持续性深度研讨》，完整统摄 计算机核心 与 领域核心：
\begin{itemize}
    \item 社会语境：5 学时
    \item 伦理分析方法：3 学时
    \item 职业伦理与专业准则：4 学时
    \item 知识产权与开源合规：2 学时
    \item 隐私与公民权利：3 学时
    \item 专业沟通与技术写作：3 学时
    \item 可持续性与绿色计算：2 学时
    \item 计算机发展简史：1 学时
    \item 计算经济学与产业生态：1 学时
    \item 安全策略、网络犯罪与法律法规：3 学时
    \item 多元、公平与包容 (DEI)：4 学时
\end{itemize}

\section{真实能力任务规范}

\begin{taskbox}[任务 1：评估部署特定 AI 驱动服务/软件/产品的伦理与社会影响]
\textbf{【任务情境】}：某企业拟将一套基于深度生成模型与算法决策的 AI 辅助招聘/信用审批服务投入商业化部署。\\
\textbf{【能力陈述】}：综合识别并评估该 AI 服务上线后将影响哪些社会利益相关群体，并精确判定其如何造成歧视性决策、隐私侵害或权益受损。\\
\textbf{【能力领域】}：理论与宏观分析 \quad \textbf{【能力单元】}：系统部署 / 方案评估 / 用户接纳 / 社会议题调适\\
\textbf{【支撑知识单元】}：人工智能/人工智能基础议题、社会伦理/社会情境与计算影响、社会伦理/伦理分析方法、社会伦理/隐私与公民自由、社会伦理/多样性、平等与包容 (DEI)\\
\textbf{【达标技能等级】}：\textbf{解释与评估} \quad \textbf{【核心要求】}：计算机核心
\end{taskbox}

\begin{taskbox}[任务 2：评估收集与使用客户/用户数据的法律与伦理意涵]
\textbf{【任务情境】}：研发团队拟在移动端 APP 和云端后端建立全埋点用户行为采集系统以支撑精准商业变现。\\
\textbf{【能力陈述】}：严密分析拟采集数据的敏感性属性，精准判定在该数据全生命周期采集、存储、聚合挖掘与流转过程中可能引发的所有法定合规风险与伦理侵权隐患。\\
\textbf{【能力领域】}：理论分析 \quad \textbf{【能力单元】}：方案评估 / 用户接纳 / 社会议题调适\\
\textbf{【支撑知识单元】}：数据管理/数据安全与隐私保护、社会伦理/伦理分析方法、社会伦理/隐私与公民自由、社会伦理/多样性、平等与包容 (DEI)\\
\textbf{【达标技能等级】}：\textbf{解释与评估} \quad \textbf{【核心要求】}：计算机核心
\end{taskbox}

\begin{taskbox}[任务 3：评估文化因素在技术产品采纳与社会公平中的决定性角色]
\textbf{【任务情境】}：某全球化跨国软件服务进军新兴发展中国家市场，遭遇用户使用习惯排斥与当地特定群体利益受损争议。\\
\textbf{【能力陈述】}：熟练运用严谨的伦理分析框架与立场理论，准确界定适用的社会文化语境，推演该技术引入对不同社会阶层、文化群体的深层伦理冲击与公平性影响。\\
\textbf{【能力领域】}：理论分析 \quad \textbf{【能力单元】}：系统部署 / 方案评估 / 社会议题调适\\
\textbf{【支撑知识单元】}：社会伦理/社会情境与计算影响、社会伦理/伦理分析方法、社会伦理/隐私与公民自由、社会伦理/多样性、平等与包容 (DEI)\\
\textbf{【达标技能等级】}：\textbf{解释与评估} \quad \textbf{【核心要求】}：计算机核心
\end{taskbox}

\begin{taskbox}[任务 4：编制企业发布技术服务所承担的问责、责任与法律责任规范文档]
\textbf{【任务情境】}：企业即将向公众正式发布一款高自主度自动驾驶系统或关键医疗辅助分析软件平台。\\
\textbf{【能力陈述】}：系统收集法务、工程质量、网络安全与行业规范信息，撰写一份条理严密、可实操的企业级责任声明与风险问责白皮书。\\
\textbf{【能力领域】}：应用与理论实践 \quad \textbf{【能力单元】}：需求工程 / 系统部署 / 文档编制\\
\textbf{【支撑知识单元】}：社会伦理/职业道德与规范、社会伦理/知识产权与开源许可、社会伦理/隐私与公民自由、社会伦理/专业沟通与表达、社会伦理/多样性、平等与包容 (DEI)\\
\textbf{【达标技能等级】}：\textbf{解释} \quad \textbf{【核心要求】}：计算机核心
\end{taskbox}

\begin{taskbox}[任务 5：制定团队跟进与企业战略相匹配的伦理、法律及职业议题策略]
\textbf{【任务情境】}：研发中心规模快速扩张，亟需建立规范机制确保工程团队在新兴技术研发中与企业伦理战略保持步调一致。\\
\textbf{【能力陈述】}：综合评估企业中长期战略目标，制定一套通过专业技术文档和日常沟通渠道向全体工程师持续宣贯、培训与追踪伦理、法律动态的行动方案。\\
\textbf{【能力领域】}：应用与工程管理 \quad \textbf{【能力单元】}：全周期管理 / 需求分析 / 方案评估 / 持续改进\\
\textbf{【支撑知识单元】}：社会伦理/专业沟通与表达、社会伦理/职业道德与规范\\
\textbf{【达标技能等级】}：\textbf{解释、评估与开发} \quad \textbf{【核心要求】}：计算机核心
\end{taskbox}

\begin{taskbox}[任务 6：将合法合规的隐私保护要求深度融入系统研发生命周期]
\textbf{【任务情境】}：架构团队正在设计全新企业级微服务分布式中台，需确保系统从代码编写、API 暴露到存储均满足 GDPR 与国家数据安全法。\\
\textbf{【能力陈述】}：识别并分解法定隐私与伦理合规指标，制定可执行的技术实施方案与团队沟通计划，确保隐私保护设计在系统开发全流程中落地。\\
\textbf{【能力领域】}：软件工程与理论实践 \quad \textbf{【能力单元】}：需求分析 / 架构设计 / 系统开发 / 流程管控 / 质量改进\\
\textbf{【支撑知识单元】}：社会伦理/职业道德与规范、社会伦理/隐私与公民自由、社会伦理/专业沟通与表达\\
\textbf{【达标技能等级】}：\textbf{解释与开发} \quad \textbf{【核心要求】}：计算机核心
\end{taskbox}

\begin{taskbox}[任务 7：在技术写作与口头汇报中展现敏锐的文化意识与文化能力]
\textbf{【任务情境】}：作为跨国开源基金会或大型跨国工程团队的负责人，需向来自不同大洲、多元文化背景的核心利益相关者发表技术演讲与撰写正式提案。\\
\textbf{【能力陈述】}：在书面技术文档与口头演说中高度展现文化敏锐度，自觉规避排他性用语与文化偏见，实现富有同理心且极具说服力的专业技术沟通。\\
\textbf{【能力领域】}：理论与人文素养 \quad \textbf{【能力单元】}：技术写作 / 社会议题调适\\
\textbf{【支撑知识单元】}：社会伦理/社会情境与计算影响、社会伦理/专业沟通与表达、社会伦理/多样性、平等与包容 (DEI)\\
\textbf{【达标技能等级】}：\textbf{解释} \quad \textbf{【核心要求】}：计算机核心
\end{taskbox}

\begin{taskbox}[任务 8：参与制定企业或研发团队的可持续发展与绿色低碳政策]
\textbf{【任务情境】}：科技企业响应国家双碳战略，计划在其全球数据中心运营与日常软件工程研发流程中推行全面碳减排。\\
\textbf{【能力陈述】}：有效将可持续计算与绿色计算知识与企业商业战略及组织使命深度融合，撰写具备行业领先水平的软件能耗优化与设备绿色循环政策。\\
\textbf{【能力领域】}：理论与工程战略 \quad \textbf{【能力单元】}：方案设计 / 系统构建 / 策略融合 / 文档编制 / 持续优化\\
\textbf{【支撑知识单元】}：社会伦理/专业沟通与表达、社会伦理/绿色计算与可持续发展\\
\textbf{【达标技能等级】}：\textbf{解释、评估与开发} \quad \textbf{【核心要求】}：计算机核心
\end{taskbox}

\begin{taskbox}[任务 9：阐释研发团队与用户多元化对企业创新与系统质量的深远收益及同质化危害]
\textbf{【任务情境】}：为企业高管层撰写决策咨询报告，论证在产研部门大力推进员工背景多元化、倾听弱势用户群体诉求对规避重大商业与公关危机的重要性。\\
\textbf{【能力陈述】}：运用严谨实证数据与典型工业惨痛事故教训，向决策层专业且令人信服地传递多元化对激发团队创新活力、消除技术偏见的决定性作用。\\
\textbf{【能力领域】}：理论与组织沟通 \quad \textbf{【能力单元】}：文档编制 / 团队管理 / 社会议题调适 / 组织改进\\
\textbf{【支撑知识单元】}：社会伦理/专业沟通与表达、社会伦理/多样性、平等与包容 (DEI)\\
\textbf{【达标技能等级】}：\textbf{解释} \quad \textbf{【核心要求】}：计算机核心
\end{taskbox}

\begin{taskbox}[任务 10：协助设计并推行旨在切实促进员工与用户群体 DEI 的系统性组织变革]
\textbf{【任务情境】}：高校计算机系或大型科技部门立项开展“多元、平等与包容（DEI）”系统性革新工程。\\
\textbf{【能力陈述】}：在具体的教学或研发组织情境中创造性应用 DEI 原理，针对人才选拔、代码评审、无障碍系统优化及职场关怀，设计并推行一套长效变革机制。\\
\textbf{【能力领域】}：应用与理论实践 \quad \textbf{【能力单元】}：需求分析 / 制度设计 / 用户接纳 / 社会议题调适 / 系统性改进\\
\textbf{【支撑知识单元】}：社会伦理/专业沟通与表达、社会伦理/多样性、平等与包容 (DEI)\\
\textbf{【达标技能等级】}：\textbf{解释、应用、评估与开发} \quad \textbf{【核心要求】}：计算机核心
\end{taskbox}

\section{工作组委员会与致谢名单}

\textbf{委员会主席：}\\
Brett A. Becker (爱尔兰都柏林大学 / University College Dublin, Ireland)

\textbf{指导委员会成员：}
\begin{itemize}
    \item Richard L. Blumenthal (美国瑞吉斯大学 / Regis University, USA)
    \item Michael Goldweber (美国泽维尔大学 / Xavier University, USA)
    \item James Prather (美国阿比林基督大学 / Abilene Christian University, USA)
    \item Susan Reiser (美国北卡罗来纳大学阿什维尔分校 / University of North Carolina Asheville, USA)
    \item Michelle Trim (美国马萨诸塞大学阿默斯特分校 / University of Massachusetts Amherst, USA)
    \item Titus Winters (美国纽约 Google 公司 / Google, Inc, USA)
\end{itemize}

\textbf{核心贡献者：}
\begin{itemize}
    \item Johanna Blumenthal (美国瑞吉斯大学 / Regis University, USA)
    \item MaryAnne Egan (美国锡耶纳学院 / Siena College, USA)
    \item Keith Quille (爱尔兰都柏林理工大学 / Technological University of Dublin, Ireland)
    \item Mehran Sahami (美国斯坦福大学 / Stanford University, USA)
    \item Mark Scanlon (爱尔兰都柏林大学 / University College Dublin, Ireland)
    \item Karren Shorofsky (美国旧金山大学法学院 / University of San Francisco School of Law, USA)
    \item Ellen Walker (美国海勒姆学院 / Hiram College, USA)
\end{itemize}

\section{附录：核心主题与技能要求等级对照表}

\begin{longtable}{p{0.20\textwidth} p{0.48\textwidth} c c c}
\toprule
\textbf{知识单元} & \textbf{核心知识点与具体主题} & \textbf{认知等级} & \textbf{分类} & \textbf{学时} \\
\midrule
\endfirsthead
\toprule
\textbf{知识单元} & \textbf{核心知识点与具体主题} & \textbf{认知等级} & \textbf{分类} & \textbf{学时} \\
\midrule
\endhead
\bottomrule
\endfoot
\multirow{2}{*}{\makecell[l]{社会伦理/社会语境}} 
& AI 演进世界的社会影响、社交媒体与算法决策融入公民生活 & 评估 & 计算机核心 & 3 \\
& 互联网演进与控制权、数字鸿沟、信息无障碍规范、情境感知计算 & 解释 & 领域核心 & 2 \\
\midrule
\multirow{2}{*}{\makecell[l]{社会伦理/伦理分析方法}} 
& 论证逻辑谬误规避、伦理学理论与技术决策框架、文化塑造作用 & 应用 & 计算机核心 & 2 \\
& 职业伦理自查清单、方案评估量规、利益相关者分析、立场理论 & 开发 & 领域核心 & 1 \\
\midrule
\multirow{2}{*}{\makecell[l]{社会伦理/职业伦理与\\专业准则}} 
& 公共价值与法律、专业素养与信托责任、学会守则、软件安全法律责任 & 评估 & 计算机核心 & 2 \\
& 公共政策参与、不良后果预警、“吹哨人”保护、跨文化道德困境、AUP 规范 & 解释 & 领域核心 & 2 \\
\midrule
\multirow{2}{*}{\makecell[l]{社会伦理/知识产权}} 
& 知识产权哲学法理基石、专有权利体系、无形数字知识产权 & 解释 & 计算机核心 & 1 \\
& DRM 技术法理、版权/专利/商密对比、代码抄袭剽窃判定、开源许可证 & 应用 & 领域核心 & 1 \\
\midrule
\multirow{2}{*}{\makecell[l]{社会伦理/隐私与公民\\权利}} 
& 事务数据库与普适监控隐私侵蚀、差分隐私数学原理、隐私增强技术 & 解释 & 计算机核心 & 2 \\
& 隐私权哲学法理渊源、重点法域隐私立法框架、表达自由 & 评估 & 领域核心 & 1 \\
\midrule
\multirow{2}{*}{\makecell[l]{社会伦理/专业沟通与\\团队协作}} 
& 源码文档提炼、API 手册编写、团队沟通规范、受众感知策略、无障碍表达 & 应用 & 计算机核心 & 2 \\
& 跨文化沟通能力、软件项目多维竞争价值观权衡、职场伦理倡导写作 & 应用 & 领域核心 & 1 \\
\midrule
\multirow{2}{*}{\makecell[l]{社会伦理/可持续性与\\绿色计算}} 
& 可持续工程实践、计算设备全生命周期生态影响、区块链共识能耗权衡 & 应用 & 计算机核心 & 1 \\
& 绿色 IT 与低碳软件设计标准、系统性生态效应、普适计算与智慧环保 & 评估 & 领域核心 & 1 \\
\midrule
\makecell[l]{社会伦理/计算机发展\\简史} 
& 机械计算前史、早期现代数字计算、微机革命、互联网与 Web、云计算与 AI & 解释 & 领域核心 & 1 \\
\midrule
\makecell[l]{社会伦理/计算经济学\\与产业生态} 
& 服务商生态与围墙花园、知识与注意力经济、人才供需、商业模式与定价、全球化分工与自动化冲击、数字金融变革 & 解释 & 领域核心 & 1 \\
\midrule
\multirow{2}{*}{\makecell[l]{社会伦理/安全策略、\\法律与网络犯罪}} 
& 计算机犯罪案例与法律追责、社会工程学欺诈、非法访问、黑客恶意软件 & 解释 & 计算机核心 & 2 \\
& 犯罪预防战略、红蓝对抗实操、编制企事业单位信息安全管理策略规范 & 应用 & 领域核心 & 1 \\
\midrule
\multirow{2}{*}{\makecell[l]{社会伦理/多元、\\公平与包容 (DEI)}} 
& 身份认同互塑、多元团队收益与同质化盲区、包容语言、无障碍开发 & 解释 & 计算机核心 & 2 \\
& 多元化行业榜样、多元化实证收益与缺乏多元化代价、技术霸权历史反思 & 评估 & 领域核心 & 2 \\
\bottomrule
\end{longtable}
